\section{ریسه در هسته‌ی ویندوز و لینوکس}

\subsection*{بخش آ – جایگاه ساختارها و \lr{TEB}}

\begin{center}\small\atlrows
\begin{tabular}{@{}c>{\centering\arraybackslash}p{2.4cm}>{\raggedleft\arraybackslash}p{7.4cm}@{}}
  \rowcolor{atl-primary}\atlhead{ساختار} & \atlhead{فضای حافظه} & \atlhead{محتوا}\\
  \lr{ETHREAD} & هسته & اشاره‌گر به پردازه و به \lr{KTHREAD}\\
  \lr{KTHREAD} & هسته & زمان‌بندی و همگام‌سازی، پشته‌ی هسته، اشاره‌گر به \lr{TEB}\\
  \lr{TEB} & کاربر & شناسه‌ی ریسه، پشته‌ی کاربر، \lr{TLS}\\
\end{tabular}
\end{center}

\textbf{چرا \lr{TEB} در فضای کاربر است؟} کد کاربر (کتابخانه‌ها و \lr{DLL}ها) مدام آن را می‌خواند و می‌نویسد (\lr{TLS}، کد آخرین خطا، حدود پشته). در فضای کاربر این کار با یک دسترسی عادی به حافظه انجام می‌شود، بدون فراخوانی سیستمی.

اگر این داده‌ها در هسته بودند، هر دسترسی به یک فراخوانی سیستمی (تغییر مد، ذخیره و بازیابی زمینه) نیاز داشت. برنامه‌های چندریسمانی که \lr{TLS} و کد خطا را زیاد می‌خوانند، به‌طور محسوس کندتر می‌شدند.

\subsection*{بخش ب – دو پشته}

\textbf{۱) چرا پشته‌ی جدا؟} پشته‌ی کاربر زیر کنترل کد کاربر است. اگر هسته روی آن کار کند:
\begin{itemize}
  \item کاربر می‌تواند \code{sp} را روی آدرس هسته، آدرس نگاشت‌نشده یا فضای ناکافی بگذارد و هسته را خراب کند یا حافظه‌اش را بنویسد؛
  \item داده‌های هسته (آدرس بازگشت، ثبات‌ها، اشاره‌گرها) برای کد کاربر و ریسه‌های دیگر پردازه خواندنی و نوشتنی می‌شود: نشت اطلاعات و دست‌کاری مسیر اجرای هسته.
\end{itemize}
پس هسته به پشته‌ی خصوصی خودش در حافظه‌ی هسته می‌رود.

\textbf{۲) اشاره‌گرها:} پشته‌ی کاربر در \lr{TEB} و پشته‌ی هسته در \lr{KTHREAD} نگهداری می‌شود.

\subsection*{بخش ج – زمان‌بندی در ویندوز}

\textbf{۱)} زمان‌بند هسته مستقیماً با \lr{KTHREAD} کار می‌کند (اولویت، وضعیت، پشته‌ی هسته، زمینه). ریسه‌ها از یک پردازه‌اند، پس فضای آدرس هم عوض نمی‌شود.

\textbf{۲)} هدف، جداسازی سازوکار از سیاست است. لایه‌ی هسته (\lr{KTHREAD}) فقط سازوکارهای پایین‌سطح و سریع دارد (زمان‌بندی، همگام‌سازی، وقفه). لایه‌ی اجرایی (\lr{ETHREAD}) مدیریت سطح‌بالا دارد (تعلق به پردازه، شناسه، امنیت). نتیجه: زمان‌بند فقط با ساختاری کوچک کار می‌کند و هر لایه مستقل از دیگری تغییر می‌کند.

\subsection*{بخش د – مدل یکپارچه‌ی لینوکس}

\textbf{۱)} زمان‌بند لینوکس همه‌ی موجودیت‌ها را یکسان و به‌صورت \lr{task} (\code{task\_struct}) می‌بیند و پردازه و ریسه را متمایز نمی‌کند. «ریسه» یعنی \lr{task}ی که با \code{clone()} ساخته شده و با \lr{task}های دیگر منابع (فضای آدرس، فایل‌ها، ...) را به اشتراک می‌گذارد. «پردازه» گروه این \lr{task}هاست (\lr{tgid} یکسان).

\textbf{۲) برنامه‌ی ۴ ریسه‌ای:}
\nopagebreak[4]
\begin{center}\small\atlrows
\begin{tabular}{@{}c>{\raggedleft\arraybackslash}p{5.8cm}>{\raggedleft\arraybackslash}p{5.8cm}@{}}
  \rowcolor{atl-primary}\atlhead{} & \atlhead{ویندوز} & \atlhead{لینوکس}\\
  ساختار پردازه & ۱ ساختار پردازه (که \lr{ETHREAD}ها به آن اشاره می‌کنند) & ندارد؛ پردازه همان گروه \lr{task}هاست\\
  ساختار ریسه & ۴ \lr{ETHREAD} با \lr{KTHREAD} خودشان & ۴ \code{task\_struct}\\
  منابع مشترک & در ساختار پردازه & ۱ مجموعه (\lr{mm}، فایل‌ها، ...) که هر ۴ \lr{task} به آن اشاره می‌کنند\\
  پشته‌ی هسته & ۴ (یکی برای هر ریسه) & ۴ (یکی برای هر \lr{task})\\
\end{tabular}
\end{center}
(۴ \lr{TEB} در فضای کاربر است، نه هسته.)

\subsection*{بخش هـ – \lr{clone} با \lr{CLONE\_VM} و \lr{CLONE\_FILES}}

\begin{center}\small\atlrows
\begin{tabular}{@{}c>{\raggedleft\arraybackslash}p{2.2cm}>{\raggedleft\arraybackslash}p{9cm}@{}}
  \rowcolor{atl-primary}\atlhead{مورد} & \atlhead{اثر در والد} & \atlhead{دلیل و پرچم}\\
  \code{open} یا \code{close(3)} & بله & \code{CLONE\_FILES}: جدول توصیفگر فایل مشترک است\\
  \code{chdir} & خیر & فرزند نسخه‌ی خودش از اطلاعات سیستم فایل (پوشه‌ی کاری) را دارد. پرچم لازم: \code{CLONE\_FS}\\
  \code{SIGINT} & خیر & فرزند جدول گردانندهای سیگنال خودش را دارد. پرچم لازم: \code{CLONE\_SIGHAND} (نیازمند \code{CLONE\_VM})\\
\end{tabular}
\end{center}

\subsection*{بخش و – انعطاف‌پذیری لینوکس در برابر ساختار سلسله‌مراتبی ویندوز}

\begin{center}\small\atlrows
\begin{tabular}{@{}c>{\raggedleft\arraybackslash}p{5.8cm}>{\raggedleft\arraybackslash}p{5.8cm}@{}}
  \rowcolor{atl-primary}\atlhead{} & \atlhead{لینوکس (پرچم‌های \lr{clone})} & \atlhead{ویندوز (ساختارهای لایه‌ای)}\\
  ایده & یک ساختار و پرچم‌هایی که تعیین می‌کنند چه چیزی مشترک باشد & ساختارهای ثابت پردازه و ریسه در لایه‌های جدا\\
  مزیت & هر ترکیب اشتراک، با یک سازوکار ساده: \lr{fork}، ریسه، و ساخت کانتینر (با پرچم‌های \lr{namespace}) & جداسازی سیاست و سازوکار؛ مرز روشن امنیتی و معنای ثابت و قابل‌پیش‌بینی\\
  هزینه & معنای «پردازه» و «ریسه» مبهم‌تر؛ خطا در انتخاب پرچم‌ها & انعطاف کم؛ مثلاً کانتینر باید روی همان ساختار ثابت ساخته شود\\
\end{tabular}
\end{center}
